\section{相关研究与设计问题}
\label{sec:related-work}

\subsection{差分传递函数与阻抗平衡}
差分测量既要求保留有效的差模响应，也要求减小共模转差模响应。Pall\`as-Areny 和 Webster 建立了具有有限差模增益与共模增益的差分放大器模型，并通过实验验证所预测的幅值和相位\citep{pallasareny1991cmrr}。其配套分析以四个传递函数描述级联差分级，并明确假设后级不会对前级产生负载作用\citep{pallasareny1991cascaded}。对于直接连接接收器的无源梯形网络，接收器阻抗会进入支路传递函数，因此需要从建立匹配方程时就纳入负载。

复增益保留了比较共模抵消条件所需的信息。Koide 利用差模增益与共模增益的复数比，解释了有意引入不平衡如何补偿有源器件的共模误差\citep{koide1997redef}。Gasulla 等指出，描述全差分电路的瞬态行为和稳定性需要同时考虑共模增益与共模转差模增益\citep{gasulla2002cmresp}。Hornero 等进一步分析了支路失配引起的频率相关共模响应，包括带通行为和相移\citep{hornero2021cm}。这些复响应分析说明，需要分别报告各增益随频率变化的复响应；仅凭抑制比的幅值无法确定哪些支路系数必须相等。

\subsection{拓扑、端接与元件匹配}
Casas 和 Pall\`as-Areny 分析了非耦合与耦合模拟差分滤波器，并给出了高阶扩展\citep{casas1996basics}。非耦合支路通过传递函数匹配实现抵消；改变互联方式还可利用电路拓扑实现抵消。其中，浮置公共节点结构为基于拓扑的共模抑制提供了重要参照。图示二阶实例由两个对地梯形网络组成，其输出端可连接等值电阻负载，并在两个输出之间接入跨接电容。因此，元件条件需要针对这些具体连接关系推导。

阻抗平衡也决定实际平衡接口的响应。Whitlock 和 Floru 展示了信号源与接收器的共模阻抗不平衡如何相互作用，并介绍了一种具有高共模输入阻抗的电路\citep{whitlock2004balancedinput}。对于本文梯形网络，需要回答的问题是：两条开路电压传递函数相同的支路，在接入负载后是否仍然相同？开路传递函数确定电压比，输出阻抗确定所接负载如何改变这一电压比。保留这两个函数，便可显式描述端接的影响，直接比较开路匹配与指定有限负载下的匹配。

\subsection{网络实现与容差分析}
经典综合的目标是由给定传递函数构造一个网络。Guillemin 建立了连分式梯形综合，并给出端接后的传递结构 $Y_{12}/(1+Y_{22})$\citep{guillemin1949rcsynthesis}。Fialkow 和 Gerst 刻画了梯形网络与一般双端口 RC 网络的传递函数及增益常数的可实现范围\citep{fialkow1951rcladder,fialkow1952rcnetworks}；其连续实现族构造还表明，同一传递函数允许许多正值实现。梯形网络的最大可实现增益由 Paige 和 Kuh 给出，Wang 等借助 Ali 和 Kim 的同余变换证明该范围内任意指定增益均可实现\citep{paige1960maximum,wang2017rcladder,ali1975continuous}。端接在这一理论中以影响可实现增益范围的方式进入\citep{fialkow1951rcladder,fialkow1952rcnetworks}。

本文提出的匹配问题与此不同：本文计算在指定端接和测量口径下带载传递函数等于固定参考支路的正值元件族。传输矩阵与连分式表示是完成该计算的已有工具\citep{guillemin1949rcsynthesis,fialkow1951rcladder,lee1978chain}，本文用其表达已知端接导纳施加的约束。所得计数区分包含绝对增益的匹配与直流校准后的匹配，不属于新的可实现性定理。

元件实现需要与灵敏度一起考虑。Moschytz 讨论了有源滤波器设计中灵敏度、元件值跨度与元件容差之间的权衡\citep{moschytz0radeffs}。Director 和 Rohrer 建立了用于计算网络灵敏度的广义伴随网络框架\citep{director1969the}。Endres 和 Kirkpatrick 研究了差分环路滤波器失配如何在快速建立的锁相环中引入额外瞬态\citep{endres0ensitivi}。这些灵敏度与瞬态研究支持将精确标称抵消条件、元件扰动影响和建立时间规范分别考察；各问题均需要明确的响应函数与评价准则。

\subsection{本文工作的范围}
本文推导任意有限级数正值对地 RC 梯形网络带载传递函数相等的系数条件。经典梯形提取提供系数坐标，多项式核给出已知有理负载下的局部匹配维数。分析进一步将原始抵消条件与校准后的动态比值测量联系起来，解析构造通过原生交流与瞬态电路输出核验。

上述抵消条件针对确定性的线性传递函数。Serrano-Finetti 等表明，共模电子噪声分析除确定性共模抑制比外，还需要噪声源及其互相关模型\citep{serranofinetti2019cmnoise}。因此，评价噪声性能需要给出相关噪声源的位置与相关性。第~\ref{sec:noise-pairing} 节在指定温度和频带下计算电阻热噪声，并保留比值读出中的通道协方差项。
